\begin{abstract}
Issue trackers contain execution histories that may support early identification of work unlikely to finish within a sprint. However, retrospective evaluations can confuse forecasting unfinished work with recognizing work already completed, and percentage-based alert budgets can behave unexpectedly in small cohorts. We conduct an archival study using TAWOS v1.1, reconstructing 28,746 issue--sprint commitments from timestamped membership changes. Of these, 22,147 have assessable outcomes, yielding 88,588 snapshots at four sprint landmarks. Fourteen projects, containing 18,322 labeled instances, meet the evaluation criteria. We compare six methods using chronological within-project splits and leave-one-project-out evaluation, with validation-only probability recalibration. At the midpoint, dynamic CatBoost increases sprint-macro recall from 44.46\% to 53.36\% over matched static CatBoost on temporal tests, a difference of 8.90 percentage points (95\% paired-bootstrap interval: 6.38--11.70). However, the difference falls to 2.59 points when evaluating only issues still open. Rounding a nominal 20\% budget upward produces a temporal mean sprint-level allocation of 41.54\%; strict downward rounding reduces dynamic recall to 18.20\%, while retaining a positive difference over the static model. Recalibration improves temporal Brier score but not cross-project Brier score. A simple execution-state rule is competitive, and its temporal recall difference from dynamic CatBoost is inconclusive. These findings support the predictive value of recorded execution histories while showing that eligibility, integer capacity, strong baselines, and calibration transfer materially affect practical conclusions. Outcomes are tracker-state proxies under archived sprint-date assumptions, not independently verified delivery or causal intervention outcomes.

\end{abstract}
\noindent\textbf{Keywords:} empirical software engineering; sprint outcome prediction; issue tracking; predictive process monitoring; probability calibration; budget-constrained alerting.
